\documentclass[12pt]{article}
\usepackage{amsmath,amssymb}
\usepackage{amsthm}
\usepackage{hyperref}
\usepackage{geometry}
\usepackage{booktabs}
\geometry{margin=1in}
\setlength{\emergencystretch}{2em} % lets TeX loosen a line rather than overrun the margin

\title{A Curvature Identity for the Fibonacci Digital-Root Cycle Modulo~9}

\author{Guy Rinzema\\
Independent researcher\\
\texttt{guy.rinzema@icloud.com}}

\date{October 2026 (Version 5)}

%--- theorem environments
\newtheorem{thm}{Theorem}
\newtheorem{prop}{Proposition}
\theoremstyle{remark}
\newtheorem{rem}{Remark}

\begin{document}
\maketitle

\begin{abstract}
For each seed we form the $24$-periodic orbit of the Fibonacci recurrence modulo~$9$,
write its terms with the digital-root convention $0\mapsto 9$, and measure its
\emph{digital curvature} $S=\sum_n \kappa_n^2$, the sum of squared discrete second
differences. On the $72$ full-period seeds modulo~$9$ the curvature takes exactly three
values, $\{648,972,1296\}$, each on $24$ seeds; the maximum $1296$ is attained precisely by
the classical Fibonacci cycle $(1,1)$ and its rotations. This is a finite statement and we
prove it by exhaustive enumeration. The curvature is the squared norm of a circulant
operator and therefore has an exact Fourier form,
$S=\frac{16}{N}\sum_{k}\sin^4(\pi k/N)\,|\hat d_k|^2$, which makes it a high-pass energy.
Modulo~$9$ the Fibonacci class is also the class with the largest share of its curvature
in the near-Nyquist half of the spectrum (about $98\%$), but this describes the maximizer
rather than explaining it: at other moduli neither the near-Nyquist energy nor its share
reliably identifies the maximizing orbit. We also show that modulo~$9$ the rule
$0\mapsto 9$ does not affect $S$ at all---raw residues $0,\dots,8$ give the same three
values---because every orbit negates itself after half a period. Finally we report how far the maximality
property extends to other moduli (it is a small-modulus phenomenon, not a law) and retire
a conjecture from earlier versions of this note. This version corrects errors in
Versions~1--4; see the revision history in Section~\ref{sec:errata}.
\end{abstract}

\noindent\textbf{Keywords:} Fibonacci numbers; Pisano period; digital root;
discrete curvature; circulant operator; discrete Fourier transform.

\section{Introduction}

Let $F(n)$ be the Fibonacci numbers, $F(0)=0$, $F(1)=1$. The sequence $F(n)\bmod 9$ has
Pisano period $24$~\cite{Wall1960}. Replacing each residue $0$ by $9$ (the digital-root
convention used in base~$10$, since the digital root of a positive integer equals its
value modulo~$9$ with $0\mapsto 9$; the digital root of $0$ itself is $0$~\cite{OEIS})
yields the familiar cycle
\[
1,1,2,3,5,8,4,3,7,1,8,9,8,8,7,6,4,1,5,6,2,8,1,9.
\]
This $24$-tuple is well known, but its higher-order difference statistics are not usually
examined. We attach to every modulo-$9$ orbit a single number---its \emph{digital
curvature}, the energy of its discrete second difference---and ask what values that number
takes.

The answer modulo~$9$ is a clean trichotomy (Theorem~\ref{thm:trichotomy}). The
curvature functional is the squared norm of the cyclic second-difference operator, which
is circulant and hence diagonalized by the discrete Fourier transform. This gives an exact
spectral formula for $S$ (Proposition~\ref{prop:spectral}), which turns the opaque integer
$S$ into an interpretable high-pass energy. The formula does not by itself prove the
classification, which rests on enumeration, and at general moduli it does not identify
the maximizing orbit (Section~\ref{sec:maximality}); we use it as a description, not a
derivation. The spectral identity is a standard fact of harmonic analysis on the cycle;
what is (to our knowledge) not previously recorded is its application to the digital-root
orbits of a Pisano period.

We have tried to be careful about what is and is not established. The trichotomy is a
theorem proved by finite computation. The spectral identity and the half-period negation
property (Proposition~\ref{prop:halfneg}) are theorems. The statement that the Fibonacci
orbit maximizes curvature is true for small moduli but fails for larger ones, and we
report it as the bounded empirical fact that it is. Earlier versions of this note advanced
a closed-form conjecture and a symmetry argument that turn out to be, respectively, a
numerical coincidence and an error; Version~4 in turn overstated what the spectral picture
explains and what had been verified. All of these are corrected below.

\section{Definitions}\label{sec:defs}

Fix a modulus $m\ge 2$ (we take $m=9$ unless stated otherwise) and let $P=\pi(m)$ be its
Pisano period~\cite{Wall1960,Koshy2018}. For a seed $(a_0,a_1)\in(\mathbb Z/m\mathbb Z)^2$ generate the orbit
\[
a_{n+1}=a_n+a_{n-1}\pmod m .
\]
We restrict attention to seeds whose orbit has \emph{minimal period exactly $P$}
(``full-period'' seeds); these are the orbits that realize the full Pisano cycle. The
\emph{rotation class} of such a seed is the set of $P$ pairs $(a_n,a_{n+1})$ occurring
along its orbit---equivalently, its $P$ cyclic shifts. Display
each residue with
\[
d(r)=\begin{cases}m & r\equiv 0,\\ r & r\in\{1,\dots,m-1\},\end{cases}
\]
so displayed values lie in $\{1,\dots,m\}$. For $m=9$ this is the base-$10$ digital root;
for other $m$ it is simply a choice of integer representatives. Write $r_n\in\{0,\dots,m-1\}$
for the raw residues, so $d_n=d(r_n)$, and set
\[
\kappa_n=d_{n+1}-2\,d_n+d_{n-1}\in\mathbb Z,
\qquad
S=\sum_{n=0}^{P-1}\kappa_n^{2},
\]
all indices taken cyclically modulo $P$. We call $S$ the \emph{digital curvature} of the
orbit. Two remarks and one proposition are worth making at the outset.

\begin{rem}[Choice of representatives]\label{rem:reps}
In general $S$ depends on which integer representatives of the residues are used. For
example, the orbit of $(1,1)$ modulo~$11$ has $S=638$ with representatives in
$\{0,\dots,10\}$ but $S=1430$ with the display $\{1,\dots,11\}$. Modulo~$9$, however,
these two conventions give the same $S$ on every orbit (Proposition~\ref{prop:halfneg}),
while the symmetric representatives $\{-4,\dots,4\}$ give yet another value ($648$ for the
$(1,1)$ orbit). So $S$ is an invariant of a chosen integer presentation of the orbit, not
of the abstract residue sequence. We fix the digital-root presentation throughout.
\end{rem}

\begin{prop}[Half-period negation]\label{prop:halfneg}
Let $M=\begin{pmatrix}0&1\\1&1\end{pmatrix}$, the matrix that sends each pair
$(a_n,a_{n+1})^{\mathsf T}$ of an orbit to the next pair $(a_{n+1},a_{n+2})^{\mathsf T}$.
Suppose $P$ is even and $M^{P/2}\equiv -I\pmod m$. Then on
every orbit
\[
d_n=m-r_{n+P/2}\qquad\text{for all }n,
\]
so $(Dd)_n=-(Dr)_{n+P/2}$ and the displayed and raw orbits have the same curvature $S$.
Their curvature spectra agree mode by mode: $|\widehat{Dd}_k|=|\widehat{Dr}_k|$ for every
$k$. Modulo~$9$ the hypothesis holds, since
$M^{12}=\begin{pmatrix}F_{11}&F_{12}\\F_{12}&F_{13}\end{pmatrix}
=\begin{pmatrix}89&144\\144&233\end{pmatrix}\equiv -I\pmod 9$.
\end{prop}

\begin{proof}
$M^{P/2}\equiv -I$ gives $r_{n+P/2}\equiv -r_n\pmod m$. If $r_n\ne 0$ then
$r_{n+P/2}=m-r_n$, so $m-r_{n+P/2}=r_n=d_n$; if $r_n=0$ then $r_{n+P/2}=0$, so
$m-r_{n+P/2}=m=d_n$. The second difference annihilates the constant $m$, so
$(Dd)_n=-(Dr)_{n+P/2}$. A sign and a cyclic shift change neither the sum of squares nor
the modulus of any Fourier coefficient.
\end{proof}

Among the $118$ moduli $m\le 120$ with even Pisano period, $M^{P/2}\equiv -I$ holds for
$43$, and the $(1,1)$ orbit has equal displayed and raw curvature at exactly those $43$.
For every $m\le 300$ with even Pisano period we found no case of equal curvature without
$M^{P/2}\equiv -I$, but we have not proved the converse.

\begin{rem}[Curvature is a cyclic energy]
$\kappa_n=(D d)_n$ where $D$ is the cyclic second-difference operator, so
$S=\lVert Dd\rVert^2$ is the discrete Dirichlet-type energy of the displayed orbit on the
cycle $C_P$. This observation drives everything in Section~\ref{sec:spectral}.
\end{rem}

\section{The Trichotomy Modulo 9}\label{sec:trichotomy}

\begin{thm}\label{thm:trichotomy}
Among the $72$ full-period seeds modulo~$9$, the digital curvature takes exactly the three
values
\[
S\in\{648,\;972,\;1296\},
\]
each attained by exactly $24$ seeds (equivalently, by one full-period orbit apiece).
Moreover $S=1296$ if and only if the seed lies in the rotation class of $(1,1)$, and then
\[
S=1296=\sum_{n=0}^{23}\kappa_n^2 .
\]
\end{thm}

\begin{proof}
There are $81$ ordered seeds $(a_0,a_1)\in\{1,\dots,9\}^2$; of these, $72$ generate orbits
of minimal period $24$ and $9$ do not. Evaluating $S$ on each full-period orbit is a
finite computation, carried out in the script of Section~\ref{sec:verify}; it returns the
multiset $\{648{:}24,\,972{:}24,\,1296{:}24\}$, with the value $1296$ occurring exactly on
the $24$ rotations of the $(1,1)$ orbit. A finite exhaustive check is a complete proof of a
statement quantified over a finite set.
\end{proof}

The split of the $81$ seeds into $72$ of period $24$, $8$ of period $8$ (both entries
divisible by $3$) and the fixed seed $(9,9)$ is already recorded in the comments to
OEIS~A030132~\cite{OEIS}; what is new here is the curvature. By
Proposition~\ref{prop:halfneg}, the same three values arise if the orbits are written with
residues in $\{0,\dots,8\}$ instead.

We stress that the enumeration \emph{is} the proof. Earlier versions of this note tried to
derive the three values from a congruence argument; that argument is vacuous (all three
values are divisible by $9$, so reduction modulo~$9$ cannot distinguish them) and has been
removed. What survives from it is a correct structural observation about \emph{which}
orbits carry the three values.

\begin{rem}[The three values are the three unit classes]\label{rem:units}
The unit group $(\mathbb Z/9\mathbb Z)^\times=\{1,2,4,5,7,8\}$ acts on orbits by
$a_n\mapsto g\,a_n$. \emph{At the level of residues}, $F_{n+12}\equiv -F_n\pmod 9$, so
negation $g=-1\equiv 8$ maps the $(1,1)$ orbit to its own shift by $12$; hence the subgroup
$\{\pm1\}$ fixes each orbit and the six units collapse to the three classes
\[
\{\pm1\}=\{1,8\},\quad \{\pm2\}=\{2,7\},\quad \{\pm4\}=\{4,5\},
\]
i.e.\ $(\mathbb Z/9\mathbb Z)^\times/\{\pm1\}$. Scaling the $(1,1)$ orbit by a
representative of each class gives precisely the three curvature values:
\[
g\in\{1,8\}\!\to\!1296,\qquad g\in\{2,7\}\!\to\!972,\qquad g\in\{4,5\}\!\to\!648,
\]
verified directly. So the trichotomy is the statement that curvature is an invariant of
$(\mathbb Z/9\mathbb Z)^\times/\{\pm1\}$, taking distinct values on its three classes.
(This residue-level symmetry is what makes the unit-class action well defined. It does
not survive as an exact antisymmetry of the second differences once the residues are
written as integers in $\{1,\dots,9\}$, or in $\{0,\dots,8\}$; see
Section~\ref{sec:halves}.)
\end{rem}

Table~\ref{tab:kappa} records the second differences of the maximal $(1,1)$ orbit, for use
below. One checks directly that $\sum_n\kappa_n=0$ and $\sum_n\kappa_n^2=1296$.

\begin{table}[!ht]
\centering
\[
\begin{array}{c|cccccccccccc}
n & 0 & 1 & 2 & 3 & 4 & 5 & 6 & 7 & 8 & 9 & 10 & 11\\\hline
a_n & 1 & 1 & 2 & 3 & 5 & 8 & 4 & 3 & 7 & 1 & 8 & 9\\
\kappa_n & 8 & 1 & 0 & 1 & 1 & -7 & 3 & 5 & -10 & 13 & -6 & -2
\end{array}
\]
\[
\begin{array}{c|cccccccccccc}
n & 12 & 13 & 14 & 15 & 16 & 17 & 18 & 19 & 20 & 21 & 22 & 23\\\hline
a_n & 8 & 8 & 7 & 6 & 4 & 1 & 5 & 6 & 2 & 8 & 1 & 9\\
\kappa_n & 1 & -1 & 0 & -1 & -1 & 7 & -3 & -5 & 10 & -13 & 15 & -16
\end{array}
\]
\caption{Second differences for the $(1,1)$ orbit. $\sum\kappa_n=0$, $\sum\kappa_n^2=1296$.}
\label{tab:kappa}
\end{table}

\section{Curvature as Spectral Energy}\label{sec:spectral}

The cyclic second-difference operator $D$ on length-$N$ sequences is circulant, hence
diagonalized by the discrete Fourier transform~\cite{Davis1979}. This yields an exact formula for $S$.

\begin{prop}\label{prop:spectral}
Let $d=(d_0,\dots,d_{N-1})$ be the displayed orbit (of length $N=P$), with unnormalized DFT
$\hat d_k=\sum_{n}d_n e^{-2\pi i kn/N}$. Then
\[
S=\frac1N\sum_{k=0}^{N-1}\mu_k^2\,|\hat d_k|^2,
\qquad
\mu_k=2\cos\!\frac{2\pi k}{N}-2=-4\sin^2\!\frac{\pi k}{N},
\]
equivalently
\[
\boxed{\;S=\frac{16}{N}\sum_{k=0}^{N-1}\sin^4\!\Big(\frac{\pi k}{N}\Big)\,|\hat d_k|^2.\;}
\]
\end{prop}

\begin{proof}
$\kappa_n=(Dd)_n=d_{n+1}-2d_n+d_{n-1}$ with indices modulo $N$, so $D$ is the circulant
matrix with symbol $e^{2\pi i k/N}-2+e^{-2\pi i k/N}=2\cos(2\pi k/N)-2=\mu_k$. The Fourier
modes $e^{2\pi i kn/N}$ are its eigenvectors with eigenvalues $\mu_k$, so
$\widehat{Dd}_k=\mu_k\hat d_k$. By Parseval, $S=\lVert Dd\rVert^2=\frac1N\sum_k|\mu_k\hat
d_k|^2=\frac1N\sum_k\mu_k^2|\hat d_k|^2$, and $\mu_k^2=16\sin^4(\pi k/N)$.
\end{proof}

This is standard harmonic analysis on the cycle; we record it only because it converts the
opaque integer $S$ into something interpretable. The weight $16\sin^4(\pi k/N)$ vanishes at
the constant mode $k=0$, grows monotonically across $k=1,\dots,\lfloor N/2\rfloor$, and is
maximal at the Nyquist mode $k=N/2$ (which exists because $N=\pi(m)$ is even for
$m\ge 3$). Thus \emph{digital curvature is a high-pass energy}: it ignores the orbit's mean
and slow variation and rewards rapid, near-alternating oscillation. Call the modes with
$\sin^2(\pi k/N)\ge\frac12$---equivalently $N/4\le k\le 3N/4$---the \emph{near-Nyquist
half} of the spectrum, and write
\[
B=\frac{16}{N}\sum_{N/4\le k\le 3N/4}\sin^4\!\Big(\frac{\pi k}{N}\Big)\,|\hat d_k|^2
\]
for the part of $S$ it carries. We call $B$ the \emph{near-Nyquist energy} and $B/S$ the
\emph{near-Nyquist share}; the two can rank orbits differently
(Section~\ref{sec:maximality}).

For the maximal $(1,1)$ orbit ($N=24$, $S=1296$), Table~\ref{tab:spectrum} gives the
per-mode contribution $c_k=\frac1N\mu_k^2|\hat d_k|^2$. Conjugate modes $k$ and $N-k$
contribute equally, so $S=c_0+2\sum_{k=1}^{11}c_k+c_{12}$. The near-Nyquist half carries
about $97.6\%$ of $S$ for this orbit; the constant and low-frequency modes carry the
remaining $2.4\%$.

\begin{table}[!ht]
\centering
\begin{tabular}{rrrr}
\toprule
$k$ & $16\sin^4(\pi k/24)$ & $|\hat d_k|^2/N$ & $c_k$\\
\midrule
0  & 0.0000  & 570.375 & 0.000\\
1  & 0.0046  & 17.919  & 0.083\\
2  & 0.0718  & 3.375   & 0.242\\
3  & 0.3431  & 19.125  & 6.563\\
4  & 1.0000  & 3.375   & 3.375\\
5  & 2.1974  & 2.331   & 5.122\\
6  & 4.0000  & 3.375   & 13.500\\
7  & 6.3385  & 2.331   & 14.774\\
8  & 9.0000  & 3.375   & 30.375\\
9  & 11.6569 & 19.125  & 222.937\\
10 & 13.9282 & 3.375   & 47.008\\
11 & 15.4595 & 17.919  & 277.022\\
12 & 16.0000 & 3.375   & 54.000\\
\bottomrule
\end{tabular}
\caption{Spectral decomposition of the $(1,1)$ orbit's curvature. Conjugate modes $k$ and
$24-k$ contribute equally, so $S=c_0+2\sum_{k=1}^{11}c_k+c_{12}=1296$. The weight
$16\sin^4(\pi k/24)$ rises with frequency; the contributions $c_k$ do not (they drop from
$k=3$ to $4$, $9$ to $10$, and $11$ to $12$). The dominant pairs are $k=11,13$ ($277.0$
each) and $k=9,15$ ($222.9$ each), with the Nyquist mode $k=12$ contributing $54.0$. The
near-Nyquist half
($\sin^2(\pi k/24)\ge\frac12$, i.e.\ $6\le k\le 18$) accounts for about $97.6\%$ of $S$.}
\label{tab:spectrum}
\end{table}

\begin{rem}[The maximizer is the highest-frequency class]
The three modulo-$9$ classes are ordered by how much of their curvature sits in the
near-Nyquist half: it is $97.6\%$ of $S$ for the $1296$ class, against $94.0\%$ and
$86.6\%$ for the $972$ and $648$ classes. So among the three full-period orbit classes
modulo~$9$, the Fibonacci class is the one whose curvature is most concentrated near the
Nyquist frequency. This is a diagnosis, not a derivation: because $S$ \emph{is} the
high-pass energy, the statement reformulates maximality spectrally but does not explain
why the $(1,1)$ seed produces the most near-Nyquist spectrum in the first place
(see Section~\ref{sec:open}). At other moduli the description does not even identify the
maximizer (Section~\ref{sec:maximality}).
\end{rem}

\subsection{The half-period asymmetry, corrected}\label{sec:halves}

The identity $F_{n+12}\equiv -F_n\pmod 9$, a case of Proposition~\ref{prop:halfneg},
suggests the second differences should satisfy $\kappa_{n+12}=-\kappa_n$, which would force the two
half-period curvature sums to be equal. \emph{This is false}, and Version~2 of this note
relied on it in error. The relation $\kappa_{n+12}=-\kappa_n$ holds for $n=1,\dots,9$ but
breaks at $n=0,10,11$---exactly the index pairs adjacent to the zeros of the orbit at
positions $11$ and $23$. Consequently the halves are \emph{unequal}:
\[
\sum_{n=0}^{11}\kappa_n^2=459,\qquad \sum_{n=12}^{23}\kappa_n^2=837,\qquad 459+837=1296,
\]
and the entire imbalance $837-459=378$ comes from those three index pairs.

Version~4 attributed the break to the display rule $d(0)=9$. That is not the cause. The
residues do satisfy $r_{n+12}\equiv -r_n$, but integer representatives turn this into an
exact antisymmetry only if the representative system is closed under negation. Neither
$\{1,\dots,9\}$ nor $\{0,\dots,8\}$ is: with either one, the representatives of $x$ and
$-x$ sum to $9$ when $x\ne 0$ but to $18$ (respectively $0$) when $x=0$. Away from the
zeros the constant offset $9$ is invisible to $D$; next to them it is not. Raw residues
therefore break the antisymmetry at the same three indices, with the two halves exchanged
($837$ and $459$, as Proposition~\ref{prop:halfneg} predicts). With the symmetric
representatives $\{-4,\dots,4\}$, which are closed under negation, the antisymmetry is
exact and the halves are equal, $324$ and $324$, for a total of $648$. In the digital-root
presentation the curvature does not split evenly across the period, and the factorization
$1296=2\cdot 648$ is not explained by half-period symmetry.

\section{How Far Maximality Extends}\label{sec:maximality}

It is natural to ask whether the Fibonacci seed maximizes digital curvature for every
modulus. It does not. Define, for any modulus $m$ with even Pisano period, $S_{\max}(m)$ as
the maximum of $S$ over full-period orbits, computed with the same display convention.

\begin{description}
\item[Small moduli.] Among moduli $m\le 14$ with even Pisano period, $(1,1)$ attains
$S_{\max}(m)$ in every case; the first failure is $m=15$.
\item[Erosion with scale.] Over the $118$ moduli $m\le 120$ with even Pisano period,
$(1,1)$ attains $S_{\max}(m)$ for $78$ of them and fails for $40$. The first three failures
are $m=15,33,40$ (and these are the only failures below $41$); the full failure list is
\[
\begin{aligned}
&15,33,40,45,49,52,53,56,59,60,64,65,67,68,71,74,80,81,83,85,\\
&86,88,90,91,92,93,97,98,103,104,105,106,107,108,109,111,115,117,118,120 .
\end{aligned}
\]
The proportion of failures broadly increases with $m$ (reaching $60$--$70\%$ among moduli
in the $80$s--$100$s), though not monotonically.
\item[No arithmetic predictor.] No simple arithmetic rule we tested separates the failures
from the successes. In particular ``$5\mid m$'' fails badly in both directions ($33$ is a
failure with $5\nmid 33$; $20,25,30,35$ are multiples of $5$ for which $(1,1)$ still wins).
\end{description}

\paragraph{What the spectral picture does not say.} Version~4 claimed that ``in every
exceptional case we examined'' the actual maximizer is the orbit with the most near-Nyquist
energy.
That is false under either reading of ``most'' (Section~\ref{sec:spectral} defines the
near-Nyquist energy $B$ and share $B/S$).
\begin{itemize}
\item \emph{Share.} At $m=40$ the maximizer, the orbit of $(0,3)$ with $S=58320$, has
share $93.7\%$, against $88.8\%$ for $(1,1)$ with $S=51920$. But the orbit of $(0,13)$ has
the same curvature as $(1,1)$, $S=51920$, and share $96.1\%$. Among the $40$ failures
listed above, no curvature maximizer has the largest share in $24$ cases, and at $m=118$
the maximizer's share is slightly below that of $(1,1)$ ($93.94\%$ against $93.99\%$).
\item \emph{Energy.} At $m=71$ the maximizer, the orbit of $(0,17)$, has $S=242110$ and
$B\approx 229974.2$, while the orbit of $(0,29)$ has the smaller $S=240406$ but the
larger $B\approx 232216.4$. Over all $118$ moduli $m\le 120$ with even Pisano period, the
curvature maximizer fails to maximize $B$ at $m=70,71,75,91,102$.
\end{itemize}
What does hold over this range is weaker: at each of the $40$ failures the maximizer has
strictly larger $B$ than the $(1,1)$ orbit. The spectral identity describes where the
curvature of an orbit lives; it does not identify which orbit has the most.

\section{What Does Not Hold}\label{sec:negatives}

We record two negative results, both retiring claims from earlier versions, because closing
off false directions is part of the contribution.

\begin{description}
\item[The closed form $S_{\max}=m(\pi(m)/2)^2$ is a coincidence.] Among the moduli
$m\le 300$ with even Pisano period it holds at exactly $m=9$ ($1296$) and $m=27$
($34992$); it also fails at $m=729$. It fails already at $m=3$ (actual $36$, predicted
$48$), and on the next powers of three at $m=81$ (actual $857304$, predicted $944784$) and
$m=243$ (actual $21677544$, predicted $25509168$). We have not proved that it fails for
every larger $m$, but nothing in the data supports it. The appealing reading
$1296=9F_{12}=9\cdot 144$ is numerology specific to $m=9$. There is no closed form for
$S_{\max}(m)$ supported by the data.
\item[The phenomenon does not carry over to $a_{n+1}=a_n+2a_{n-1}$.] This recurrence is
invertible modulo $m$ exactly when $m$ is odd. With the same display convention, for every
odd $m$ with $5\le m\le 101$ the orbit of $(0,1)$ (which contains $(1,1)$) does not attain
the largest $S$ among orbits of the same minimal period; at $m=3$ it is the only orbit of
that period. Version~4 also made claims about the recurrence $a_{n+1}=2a_n+a_{n-1}$ and the
Tribonacci recurrence without stating definitions or a search range; we withdraw them.
\end{description}

\section{Computational Verification}\label{sec:verify}

The script \texttt{papers/v5/reproduce.py} in the project repository asserts each exact
numerical claim in this note, and each finite-range claim over the range stated: the
trichotomy and seed-period split (Theorem~\ref{thm:trichotomy}); the representative
comparisons of Remark~\ref{rem:reps} and Section~\ref{sec:halves}; the half-period
negation property and its finite-range companion (Proposition~\ref{prop:halfneg}); the
unit-class scalings (Remark~\ref{rem:units}); the second differences and half sums
(Table~\ref{tab:kappa}); the spectral identity on five orbits
(Proposition~\ref{prop:spectral}); the spectral table and the three near-Nyquist shares;
the maximality sweep, full failure list and decade proportions; the near-Nyquist energy and
share counterexamples and the count of failures without a unit-multiple maximizer
(Sections~\ref{sec:maximality} and~\ref{sec:open}); the closed-form check; and the
$a_{n+1}=a_n+2a_{n-1}$ sweep (Section~\ref{sec:negatives}). Spectral quantities are
computed in floating point, and every comparison the script asserts has a margin far
larger than rounding error. It uses only the Python standard library.

The original \texttt{verify\_curvature.py}, which accompanies Versions~1--3, independently
checks the trichotomy only. The Version~5 corrections follow an independent fact-check of
Version~4, which re-derived its numerical claims with separate code.

\section{Revision History and Errata}\label{sec:errata}

This is Version~5. It supersedes Versions~1--4
(\href{https://doi.org/10.5281/zenodo.15811767}{doi:10.5281/zenodo.15811767}), which
remain available as the historical record.

\paragraph{Corrections in Version~5 (to Version~4).}
\begin{itemize}
\item \textbf{The spectral maximizer claim was false.} Version~4 said that in every
exceptional case it examined the curvature maximizer is the orbit with the most
near-Nyquist energy, and asked for a proof that this always holds. Counterexamples exist for both the energy
and the share readings, including $m=40$ and $m=71$ (Section~\ref{sec:maximality}).
\item \textbf{The convention-dependence remark was false for $m=9$.} Version~4 said that
residues in $\{0,\dots,m-1\}$ give a different $S$. Modulo~$9$ they give the same $S$ on
every orbit (Proposition~\ref{prop:halfneg}); the dependence is real at other moduli
(Remark~\ref{rem:reps}).
\item \textbf{The half-period asymmetry was misattributed.} It is caused by using
representatives not closed under negation, not by the rule $d(0)=9$; raw residues break
it at the same indices (Section~\ref{sec:halves}).
\item \textbf{The verification claim was overstated.} Version~4 said its script
reproduces every numerical claim and that every claim had been checked by two independent
implementations. Its script did not test the spectral maximizer claim, the
other-recurrence claims, or the unit-multiple count, printed the near-Nyquist percentages
without asserting them, and checked the closed form only at five moduli; the original
script checks only the trichotomy. Section~\ref{sec:verify} now
states what is checked.
\item \textbf{Finite searches were stated as universal.} ``It never holds off the powers
of three'' is replaced by the range actually searched (Section~\ref{sec:negatives}).
\item \textbf{The other-recurrence claims are made precise or withdrawn}
(Section~\ref{sec:negatives}).
\item \textbf{Wording.} The abstract and introduction said the spectral identity explains
the classification and that the maximizer follows from it; it only reformulates it. The
spectral-table caption said the contributions rise with frequency; only the weight does.
The abstract wrote $\hat a_k$ for $\hat d_k$. The digital-root remark now says
``positive integer'' (the digital root of $0$ is $0$). The assumption $m\ge 2$ is stated,
and the Nyquist mode is noted to exist only because $\pi(m)$ is even for $m\ge 3$. The
failure list is now given in full. The seed-period split is credited to OEIS~A030132.
\end{itemize}

\paragraph{Corrections in Version~4 (to Versions~1--3).}
\begin{itemize}
\item \textbf{The trichotomy is proved by enumeration.} The congruence/``scaling''
argument offered as a proof in Versions~1--3 is vacuous (all three curvature values are
divisible by~$9$) and has been replaced by the finite verification, which is the actual
proof. The valid residue of that argument---the three values as the three classes of
$(\mathbb Z/9\mathbb Z)^\times/\{\pm1\}$---is kept as Remark~\ref{rem:units}.
\item \textbf{The half-period symmetry claim was wrong.} Version~2 asserted
$S=2\sum_{n=0}^{11}\kappa_n^2$ from an antisymmetry that fails for these integer
representatives; the true halves are $459$ and $837$ (Section~\ref{sec:halves}).
Version~3 already withdrew the explicit formula.
\item \textbf{The conjecture $S_{\max}=m(\pi(m)/2)^2$ is false}
(Section~\ref{sec:negatives}); it is a two-point coincidence at $m=9,27$.
\item \textbf{Section~6's ``explanation'' was replaced} by the spectral identity
(Section~\ref{sec:spectral}). Version~4 called the identity the reason the maximum is what
it is; that overstated it (see above).
\item \textbf{The novelty claim is softened.} We claim only that we are not aware of prior
work applying this spectral/Dirichlet-energy lens to the digital-root orbits of a Pisano
period; the underlying circulant diagonalization is classical. A search of OEIS
(A030132 and neighbours) and of MathSciNet for ``digital root'' and ``Pisano period''
turned up no prior treatment of these second-difference statistics, but we make no
exhaustive-priority claim.
\end{itemize}

\section{Open Questions}\label{sec:open}

\begin{itemize}
\item Characterize the moduli for which $(1,1)$ fails to maximize curvature
(Section~\ref{sec:maximality}). No arithmetic predictor was found, and $15$ of the $40$
failures with $m\le 120$ have no maximizing orbit that is a unit multiple of the Fibonacci
orbit.
\item Find a criterion, spectral or arithmetic, that does identify the curvature
maximizer; neither the near-Nyquist energy $B$ nor the share $B/S$ does
(Section~\ref{sec:maximality}).
\item Prove or refute the converse of Proposition~\ref{prop:halfneg}: is the curvature of
the $(1,1)$ orbit the same in both presentations only when $M^{P/2}\equiv -I\pmod m$? This
holds for every $m\le 300$ with even Pisano period.
\item Is the $S_{\max}=m(\pi(m)/2)^2$ coincidence at $m=9,27$ part of any larger pattern,
or genuinely isolated?
\end{itemize}

\section*{Disclosure statement}
\noindent The author reports no competing interests.

\begin{thebibliography}{9}

\bibitem{Wall1960}
D.~D.~Wall,
``Fibonacci Series Modulo $m$,''
\textit{Amer.\ Math.\ Monthly} \textbf{67} (1960), 525--532.

\bibitem{OEIS}
OEIS~A030132, ``Digital root of Fibonacci($n$).''

\bibitem{Koshy2018}
T.~Koshy,
\textit{Fibonacci and Lucas Numbers with Applications},
Wiley (2018).

\bibitem{Davis1979}
P.~J.~Davis,
\textit{Circulant Matrices},
Wiley (1979).

\end{thebibliography}

\end{document}
